\chapter{生命週期}
\label{ch:lifetime}

第~2~章的懸置參考和第~3~章「走訪時不要 \texttt{push\_back}」，都跟物件
什麼時候建立、什麼時候消失有關。 C 的變數消失時不會執行任何程式，記憶體
還給系統就結束了。 C++ 的物件消失時會執行一段程式，\texttt{string} 和
\texttt{vector} 用這段程式把記憶體還回去，所以不用替它們 \texttt{free}。

這一章先整理物件能活多久，再看物件建立和銷毀的順序，以及物件銷毀之後還
去用它會怎樣。 後半章講 \texttt{new}／\texttt{delete}，以及平常怎麼用
RAII 和智慧指標取代它們，最後談複製和移動。

\section{三種儲存期}
\label{sec:storage}

\begin{crecap}[區域變數、static 與 malloc]
C 的變數依照宣告的位置和方式，記憶體配置和釋放的時機不同：

\begin{code}{c}{C 的三種變數}
int total = 0;               // 全域：程式開始就有，程式結束才釋放

void f(void) {
    int x = 0;               // 區域：進入區塊時配置，離開區塊就釋放
    static int calls = 0;    // static 區域：只初始化一次，一直存在
    int *p = malloc(sizeof(int));   // heap：直到 free(p) 才釋放
    /* ... */
}
\end{code}

區域變數離開區塊就沒了，所以函式不能回傳區域變數的位址；static 區域變數
一直都在，回傳它的位址沒問題。 \texttt{malloc} 拿到的記憶體不會自己消失，
忘了 \texttt{free} 就是記憶體洩漏，\texttt{free} 之後還去用就是懸置指標。
\end{crecap}

C++ 沿用同一套分類，名稱叫\textbf{儲存期}（storage duration）：

\begin{center}
\begin{tabular}{llll}
\toprule
儲存期 & 例子 & 何時建立 & 何時銷毀 \\
\midrule
自動（automatic） & 區域變數、參數 & 執行到宣告 & 離開所在的區塊 \\
靜態（static） & 全域變數 & \texttt{main} 開始之前 & 程式結束 \\
 & \texttt{static} 區域變數 & 第一次執行到宣告 & 程式結束 \\
動態（dynamic） & \texttt{new} 出來的物件 & \texttt{new} & \texttt{delete} \\
\bottomrule
\end{tabular}
\end{center}

\texttt{new} 和 \texttt{delete} 是 C++ 版的 \texttt{malloc} 和
\texttt{free}，第~\ref{sec:new-delete}~節介紹。 先看前兩種：

\begin{example}[label=ex:storage]{cpp}{(storage.cpp) 自動與靜態儲存期}
#include <iostream>
using namespace std;

int nextTicket() {
    static int ticket = 100;   // static 區域變數：只初始化一次
    ticket++;
    return ticket;
}

int main() {
    for (int i = 0; i < 3; i++) {
        int step = 10 + i;     // 自動：每一輪重新建立
        cout << nextTicket() << ' ' << step << '\n';
    }
    return 0;
}
\end{example}

\iosample{}{
101 10 \\
102 11 \\
103 12 \\
}

\texttt{step} 每一輪都是新的，上一輪的值不會留下來；\texttt{ticket} 只在
第一次呼叫 \texttt{nextTicket} 時設成 $100$，之後每次呼叫都沿用同一個，
所以號碼一直往上加。

\subsection{範圍與生命週期不是同一件事}

這兩個詞要分開：

\begin{itemize}
\item \textbf{範圍}（scope）：名字在原始碼的哪一段看得到。 這是編譯時的
概念，編譯器用它決定一個名字指的是哪個變數。
\item \textbf{生命週期}（lifetime）：物件在執行時從建立到銷毀的那段期間。
物件只有在生命週期內可以使用。
\end{itemize}

自動變數的兩者剛好重疊：\texttt{step} 的名字只在迴圈本體裡看得到，物件也
只活一輪。 其他情況兩者不一致：

\begin{itemize}
\item \texttt{ticket} 的範圍只有 \texttt{nextTicket} 的本體，生命週期卻
是整個程式。 \texttt{main} 叫不出這個名字，但它一直都在。
\item 動態配置的物件根本沒有名字，只能透過指標找到。 指標變數離開範圍時，
物件不會跟著消失；沒有其他指標記得它，它就洩漏了。
\item 反過來，名字還在範圍內，它指向的物件卻可能已經死了。 指標
\texttt{p} 指著一個區塊裡的區域變數，區塊結束後 \texttt{p} 還能用，但
\texttt{*p} 已經沒有東西，這就是第~\ref{sec:dangling}~節的懸置指標。
\end{itemize}

\subsection{常見錯誤}

\begin{itemize}
\item \textbf{以為區域變數會保留上一次的值}：把 \texttt{nextTicket} 裡的
\texttt{static} 拿掉，每次呼叫都重新建立一個 \texttt{ticket}，三次都回傳
$101$。 要跨呼叫保留值，就要加 \texttt{static}。
\end{itemize}

\section{物件何時建立、何時銷毀}
\label{sec:ctor-dtor}

C++ 的結構可以寫兩個特別的函式，在物件建立和銷毀時自動執行：

\begin{itemize}
\item \textbf{建構子}（constructor）：名稱跟結構相同，沒有回傳型別。 物件
建立時自動呼叫，通常用來設定初始值。
\item \textbf{解構子}（destructor）：名稱是 \texttt{\textasciitilde} 加上
結構名稱，沒有參數也沒有回傳型別。 物件銷毀時自動呼叫，通常用來釋放資源。
\end{itemize}

寫法的細節是第~5~章的內容。 這裡只用它們在建立和銷毀時各印一行，看物件
什麼時候建立、什麼時候銷毀。

\begin{example}[label=ex:lifetime-order]{cpp}{(lifetime\_order.cpp) 會報告出生與死亡的 Tracer}
#include <iostream>
#include <string>
using namespace std;

struct Tracer {
    string name;
    Tracer(const string &n) {      // 建構子：物件建立時自動呼叫
        name = n;
        cout << "建立 " << name << '\n';
    }
    ~Tracer() {                    // 解構子：物件銷毀時自動呼叫
        cout << "銷毀 " << name << '\n';
    }
};

void f() {
    static Tracer s("static s");
}

int main() {
    cout << "main 開始" << '\n';
    Tracer a("a");
    {
        Tracer b("b");
        Tracer c("c");
        cout << "內層區塊結束" << '\n';
    }
    f();
    f();
    Tracer("暫時物件");
    cout << "main 結束" << '\n';
    return 0;
}
\end{example}

\iosample{}{
main 開始 \\
建立 a \\
建立 b \\
建立 c \\
內層區塊結束 \\
銷毀 c \\
銷毀 b \\
建立 static s \\
建立 暫時物件 \\
銷毀 暫時物件 \\
main 結束 \\
銷毀 a \\
銷毀 static s \\
}

\texttt{Tracer a("a");} 建立一個 \texttt{Tracer}，括號裡的字串交給建構子
的參數 \texttt{n}。 對照輸出逐段看：

\begin{itemize}
\item \textbf{區塊裡的物件}依宣告順序建立，離開區塊時以\textbf{相反順序}
銷毀：\texttt{b}、\texttt{c} 建立，\texttt{c}、\texttt{b} 銷毀。 後建立
的物件可能用到先建立的物件，所以先銷毀後建立的。
\item \textbf{\texttt{static} 區域物件}在第一次呼叫 \texttt{f} 時建立，
第二次呼叫不會再建立（一般的區域物件則是每次呼叫都建立、銷毀一次）。
\item \textbf{暫時物件}：\texttt{Tracer("暫時物件");} 建立一個沒有名字的
物件，它在這個運算式結束時（大致就是到分號為止）就銷毀。
\item \texttt{main} 的 \texttt{return} 之後，\texttt{a} 先銷毀，靜態儲存
期的 \texttt{s} 在程式結束時才銷毀。 全域物件則在 \texttt{main} 開始之前
就建立。
\end{itemize}

\texttt{string}、\texttt{vector} 也有解構子，負責把自己配置的記憶體還回
去。 第~3~章說 \texttt{vector} 離開範圍時記憶體會自動還回去，就是解構子
做的。

\subsection{常見錯誤}

\begin{itemize}
\item \textbf{替建構子寫回傳型別}：寫成 \texttt{void Tracer(const string
\&n)} 會編譯失敗，\texttt{g++} 報 \texttt{return type specification for
constructor invalid}。 建構子和解構子都不寫回傳型別，連 \texttt{void} 也
不寫。
\item \textbf{以為物件在最後一次使用後就銷毀}：\texttt{a} 在
\texttt{main} 開頭之後就沒再用到，仍然活到 \texttt{main} 結束。 自動物件
的銷毀時機只看它在哪個區塊，跟最後一次用到它的位置無關。
\end{itemize}

\section{懸置指標與懸置參考}
\label{sec:dangling}

物件的生命週期結束後，原本指向它的指標或綁定它的參考不會自動失效，位址
還在，只是那塊記憶體已經不屬於這個物件。 這樣的指標叫\textbf{懸置指標}
（dangling pointer），參考叫\textbf{懸置參考}（dangling reference）。
透過它們讀寫是未定義行為：可能印出舊值、垃圾，也可能當掉。 結果不固定，
所以這一節的錯誤都用 \textbf{AddressSanitizer}（簡稱 ASan）跑一次。 它是 \texttt{g++} 內建的
檢查工具，編譯時加上 \texttt{-fsanitize=address}，程式執行時一碰到已經
釋放或已經失效的記憶體，就停下來印出報告：

\begin{code}{bash}{用 ASan 編譯並執行}
g++ -std=c++17 -Wall -Wextra -g -fsanitize=address vec_dangling.cpp -o vec_dangling
./vec_dangling
\end{code}

\texttt{-g} 讓報告裡有行號。 加了 ASan 的程式會變慢、變大，只在開發和除錯
時用。 本書的 ASan 輸出取自 Linux 上的 \texttt{g++}，只節錄重點並刪掉目錄
路徑；位址和 \texttt{==} 之間的數字每次執行都不一樣。

\subsection{\texttt{vector} 換記憶體}

第~3~章提過，\texttt{push\_back} 可能讓 \texttt{vector} 換一塊更大的記憶
體。 \texttt{v.capacity()} 是目前這塊記憶體放得下幾個元素，
\texttt{size()} 超過它的時候就要換。

\begin{example}[label=ex:vec-dangling]{cpp}{(vec\_dangling.cpp) push\_back 之後，舊的參考懸置了}
#include <iostream>
#include <vector>
using namespace std;

int main() {
    vector<int> v = {10, 20, 30};
    int &first = v[0];
    int *last = &v[2];
    cout << first << ' ' << *last << ' ' << v.capacity() << '\n';

    v.push_back(40);         // 空間不夠，換一塊更大的記憶體
    cout << v.capacity() << '\n';
    cout << first << '\n';   // first 綁著舊記憶體：未定義行為
    return 0;
}
\end{example}

\begin{code}{text}{用 ASan 編譯後執行（節錄）}
10 30 3
6
=================================================================
==481641==ERROR: AddressSanitizer: heap-use-after-free on address 0x502000000010 at pc 0x57c85395f7e5 bp 0x7ffd17c330a0 sp 0x7ffd17c33090
READ of size 4 at 0x502000000010 thread T0
    #0 0x57c85395f7e4 in main vec_dangling.cpp:13

freed by thread T0 here:
    #6 0x57c85395fcff in std::vector<int, std::allocator<int> >::push_back(int&&) /usr/include/c++/13/bits/stl_vector.h:1299
    #7 0x57c85395f76b in main vec_dangling.cpp:11
\end{code}

報告從上往下看：\texttt{heap-use-after-free} 是讀了一塊已經釋放的 heap
記憶體，\texttt{READ of size 4} 是讀了一個 \texttt{int}，發生在第 $13$ 行
的 \texttt{cout << first}；\texttt{freed by} 說這塊記憶體是第 $11$ 行的
\texttt{push\_back} 釋放的（完整的報告還有一段
\texttt{previously allocated by}，說它是第 $6$ 行建立 \texttt{v} 時配置的）。
所以經過是：\texttt{v} 原本的容量是
$3$，第四個元素放不下，\texttt{push\_back} 配置一塊容量 $6$ 的新記憶體、
把元素搬過去、釋放舊的，\texttt{first} 和 \texttt{last} 還指著舊的那塊。
容量怎麼成長由實作決定，$3$ 和 $6$ 是本書測試環境的結果。 不管在哪個環境，
\texttt{size()} 超過 \texttt{capacity()} 時，所有指向元素的參考、指標和
迭代器（第~9~章）都會失效。

修法有三種：存下標而不是參考（\texttt{v[0]} 每次都重新算位址）；在
\texttt{push\_back} 之後才取參考；或者事先用 \texttt{v.reserve(n)} 把容量
預留到至少 $n$，\texttt{size()} 不超過 $n$ 之前就不會換記憶體。

\subsection{離開區塊的區域變數}

第~\ref{sec:storage}~節提過的情況，寫成程式是這樣：

\begin{code}{cpp}{指標比它指的變數活得久（錯誤示範）}
int *p = nullptr;
{
    int x = 42;
    p = &x;
}                       // x 的生命週期在這裡結束
cout << *p << '\n';     // p 還在範圍內，*p 是未定義行為
\end{code}

ASan 回報 \texttt{stack-use-after-scope}（範圍結束後還在用 stack 上的變數），
報告裡會指出讀到的是 \texttt{'x'}。

第~2~章的「回傳區域變數的參考」是同一類錯誤，只是區塊換成了函式。
\texttt{g++ -Wall} 會先警告 \texttt{reference to local variable 'x'
returned}。 在本書測試的環境下，\texttt{g++} 讓這種函式直接回傳空位址，
ASan 報的是 \texttt{SEGV on unknown address 0x000000000000}。 修法都一樣：
改成傳值回傳。

\subsection{暫時物件的 \texttt{c\_str()}}

要把 \texttt{string} 交給只收 C 字串的函式（例如 \texttt{printf}、
\texttt{fopen}），要呼叫 \texttt{s.c\_str()}，它回傳一個
\texttt{const char *}，指向 \texttt{s} 內部的字元。 這個指標只在
\texttt{s} 還活著、而且沒被修改時有效：

\begin{code}{cpp}{指向暫時字串的指標（錯誤示範）}
string first = "Grace", last = "Hopper";
const char *p = (first + " Brewster " + last).c_str();
printf("%s\n", p);      // 暫時字串已經銷毀
\end{code}

\texttt{first + " Brewster " + last} 是一個暫時物件，第 $2$ 行的分號一到
就銷毀了，\texttt{p} 隨即懸置。 ASan 報 \texttt{heap-use-after-free}，
\texttt{freed by} 指向宣告 \texttt{p} 的那一行。 這段程式 \texttt{-Wall -Wextra} 不會
警告。

暫時物件活到運算式結束，所以在同一個運算式裡用完就沒問題：
\texttt{printf("\%s\textbackslash n", (first + " Brewster " +
last).c\_str());}。 或者先存成變數 \texttt{string full = ...;}，再用
\texttt{full.c\_str()}。

\subsection{常見錯誤}

\begin{itemize}
\item \textbf{程式沒當掉就以為沒事}：未定義行為不保證當掉。 不開 ASan 時，
範例~\ref{ex:vec-dangling} 在本書測試的環境下照樣跑完，最後一行印出一個
每次執行都不同的數字；換個環境，它也可能剛好印出 \texttt{10}，看起來沒有
問題。 懷疑有懸置的時候，加 \texttt{-fsanitize=address} 跑一次。
\item \textbf{把 \texttt{c\_str()} 的結果存起來，之後又修改字串}：
\texttt{const char *p = s.c\_str();} 之後做 \texttt{s += "..."}，
\texttt{s} 可能換了一塊記憶體，\texttt{p} 就懸置了。 要用 C 字串時再呼叫
\texttt{c\_str()}，不要把結果存著。
\end{itemize}

\section{\texttt{new} 與 \texttt{delete}}
\label{sec:new-delete}

\begin{crecap}[malloc 與 free 的規則]
\texttt{malloc(n)} 向 heap 要 $n$ 個位元組，內容是垃圾值，失敗時回傳
\texttt{NULL}；\texttt{free(p)} 把它還回去。 要遵守三條規則：每塊記憶體剛好
\texttt{free} 一次；\texttt{free} 之後不能再用；最後一個指向它的指標消失
之前要 \texttt{free}。 違反第一條是雙重釋放
（double free），第二條是懸置指標，第三條是記憶體洩漏（memory leak）。
\end{crecap}

C++ 動態配置用的是運算子 \texttt{new} 和 \texttt{delete}。 單一物件用
\texttt{new}／\texttt{delete}，例如 \texttt{int *p = new int(42);} 配置一個
初始值 $42$ 的 \texttt{int}，用完 \texttt{delete p;}；陣列用
\texttt{new[]}／\texttt{delete[]}。 下面的 \texttt{Tracer} 跟
範例~\ref{ex:lifetime-order} 的一樣，只是寫得緊湊一點：

\begin{example}[label=ex:new-delete]{cpp}{(new\_delete.cpp) 配置單一物件、陣列與結構}
#include <iostream>
#include <string>
using namespace std;

struct Tracer {
    string name;
    Tracer(const string &n) { name = n; cout << "建立 " << name << '\n'; }
    ~Tracer() { cout << "銷毀 " << name << '\n'; }
};

int main() {
    int n = 5;
    int *a = new int[n];           // n 個 int 的陣列
    for (int i = 0; i < n; i++) {
        a[i] = i * i;
    }
    cout << a[n - 1] << '\n';
    delete[] a;                    // 陣列要用 delete[]

    Tracer *t = new Tracer("heap 上的 t");
    cout << t->name << '\n';
    delete t;                      // 先呼叫解構子，再釋放記憶體
    return 0;
}
\end{example}

\iosample{}{
16 \\
建立 heap 上的 t \\
heap 上的 t \\
銷毀 heap 上的 t \\
}

跟 \texttt{malloc}／\texttt{free} 比較：

\begin{center}
\begin{tabular}{lll}
\toprule
 & \texttt{malloc}／\texttt{free} & \texttt{new}／\texttt{delete} \\
\midrule
要多少 & 位元組數 & 型別（和個數） \\
回傳型別 & \texttt{void *}，C++ 裡要轉型 & 對應型別的指標 \\
建構子、解構子 & 不呼叫 & 呼叫 \\
失敗時 & 回傳 \texttt{NULL} & 丟出例外 \texttt{std::bad\_alloc} \\
陣列 & 同一組函式 & \texttt{new[]} 配 \texttt{delete[]} \\
\bottomrule
\end{tabular}
\end{center}

差別最大的是第三列。 \texttt{Tracer} 裡有一個 \texttt{string}，用
\texttt{malloc} 拿到的只是一塊沒有初始化的位元組，裡面的 \texttt{string}
從來沒有被建構過，拿來用是未定義行為。 \texttt{new} 先配置記憶體再呼叫建
構子，\texttt{delete} 先呼叫解構子再釋放記憶體。 所以
C++ 的程式裡，含有 \texttt{string}、\texttt{vector} 這類成員的結構不能用
\texttt{malloc}。

另外，\texttt{new int[n]} 跟 \texttt{malloc} 一樣不會把元素設成 $0$；要
全部是 $0$ 就寫 \texttt{new int[n]()}。

\subsection{常見錯誤}

\texttt{malloc} 的三種錯誤 \texttt{new} 都會犯，另外還多一種：配錯對。
這些錯誤都編得過，下面附上 ASan 的報告。

\begin{itemize}
\item \textbf{\texttt{new[]} 配 \texttt{delete}}：\texttt{int *a = new
int[5];} 之後寫 \texttt{delete a;}，是未定義行為。 \texttt{g++ -Wall}
這次有警告（\texttt{called on pointer returned from a mismatched
allocation function [-Wmismatched-new-delete]}），ASan 執行時則報
\texttt{alloc-dealloc-mismatch (operator new [] vs operator delete)}。
規則是成對：\texttt{new} 配 \texttt{delete}、\texttt{new[]} 配
\texttt{delete[]}、\texttt{malloc} 配 \texttt{free}，不能交叉。
\item \textbf{重複 \texttt{delete}}：同一個指標 \texttt{delete} 兩次，ASan
報 \texttt{attempting double-free}，並列出兩次釋放各在哪一行。
\texttt{delete nullptr} 什麼都不做，所以
\texttt{delete} 之後把指標設成 \texttt{nullptr} 能擋掉一部分這類錯誤。
\item \textbf{忘了 \texttt{delete}}：最常見的情況是函式有好幾個出口，其中
一個忘了釋放：
\begin{code}{cpp}{(leak.cpp) 提早 return 造成洩漏（錯誤示範，main 省略）}
#include <iostream>
using namespace std;

// 0*0 + 1*1 + ... + (n-1)*(n-1)，超過 100 就回傳 -1
int sumSquares(int n) {
    int *a = new int[n];
    for (int i = 0; i < n; i++) {
        a[i] = i * i;
    }
    int s = 0;
    for (int i = 0; i < n; i++) {
        if (s + a[i] > 100) {
            return -1;           // 提早離開，忘了 delete[]
        }
        s += a[i];
    }
    delete[] a;
    return s;
}
\end{code}
\texttt{main} 裡只有一行 \texttt{cout << sumSquares(5) << ' ' <<
sumSquares(10) << '\textbackslash n';}，印出 \texttt{30 -1}，看不出問題。 ASan 會在程式結束時檢查還有哪些記
憶體沒還：
\begin{code}{text}{ASan 的洩漏報告（節錄）}
==482383==ERROR: LeakSanitizer: detected memory leaks

Direct leak of 40 byte(s) in 1 object(s) allocated from:
    #0 0x70254f4fe6c8 in operator new[](unsigned long) ../../../../src/libsanitizer/asan/asan_new_delete.cpp:98
    #1 0x5c27037fc2d7 in sumSquares(int) leak.cpp:6
\end{code}
$40$ 個位元組就是 \texttt{sumSquares(10)} 那次的 $10$ 個 \texttt{int}。
洩漏檢查（LeakSanitizer）在 Linux 上預設開啟；macOS 內建的編譯器通常不支
援這一項，其他檢查照樣可用。
\end{itemize}

\section{RAII 與智慧指標}
\label{sec:raii}

上一節的洩漏，修法是不要自己 \texttt{new}：把 \texttt{leak.cpp} 第
$6$ 行的 \texttt{int *a = new int[n];} 換成 \texttt{vector<int> a(n);}
（記得 \texttt{\#include <vector>}），刪掉 \texttt{delete[] a;} 那一行，
其他不用改。

\texttt{a} 現在是自動變數，函式不管從哪一個 \texttt{return} 離開，都會執
行 \texttt{a} 的解構子，把記憶體還回去。 用 ASan 跑同樣的 \texttt{main}，
輸出還是 \texttt{30 -1}，沒有報告。

這個做法叫 \textbf{RAII}（Resource Acquisition Is Initialization，取得
資源就是初始化），意思是：

\begin{itemize}
\item 資源（記憶體、開啟的檔案……）交給一個物件管。 物件的建構子取得資源，
解構子釋放資源。
\item 這個物件放在自動變數裡。 區塊結束時解構子執行，資源跟著釋放，不必
自己記得釋放。
\end{itemize}

\texttt{string} 和 \texttt{vector} 都是這樣寫的，所以本書到目前為止沒有
寫過 \texttt{delete}。 第~5~章會自己寫一個這樣的類別。

\subsection{\texttt{unique\_ptr} 與 \texttt{shared\_ptr}}

有時候需要的是 heap 上的單一物件，例如要讓物件活得比建立它的函式更久，
或者第~7~章用指標處理不同型別的物件。 這時用 \texttt{<memory>} 裡的
\textbf{智慧指標}（smart pointer），也就是套用 RAII 的指標：

\begin{itemize}
\item \texttt{unique\_ptr<T>}：\textbf{獨占}一個物件。 同一時間只有一個
\texttt{unique\_ptr} 擁有它，\texttt{unique\_ptr} 銷毀時物件跟著
\texttt{delete}。 用 \texttt{make\_unique<T>(建構子的引數)} 建立。
\item \texttt{shared\_ptr<T>}：好幾個 \texttt{shared\_ptr} \textbf{共同
擁有}一個物件。 它們一起記著目前有幾個擁有者（\texttt{use\_count()}），
最後一個擁有者銷毀或放手時，物件才 \texttt{delete}。 用
\texttt{make\_shared<T>(...)} 建立。
\end{itemize}

兩者用起來都像一般指標：\texttt{*p} 取物件，\texttt{p->name} 取成員，可
以跟 \texttt{nullptr} 比較。

下面的範例沿用同一個 \texttt{Tracer}：

\begin{example}[label=ex:smart-ptr]{cpp}{(smart\_ptr.cpp) unique\_ptr 交出所有權，shared\_ptr 計算擁有者}
#include <iostream>
#include <memory>
#include <string>
using namespace std;

struct Tracer {
    string name;
    Tracer(const string &n) { name = n; cout << "建立 " << name << '\n'; }
    ~Tracer() { cout << "銷毀 " << name << '\n'; }
};

unique_ptr<Tracer> makeTracer(const string &name) {
    unique_ptr<Tracer> p = make_unique<Tracer>(name);
    cout << "makeTracer 結束" << '\n';
    return p;                      // 所有權交給呼叫者
}

int main() {
    unique_ptr<Tracer> u = makeTracer("u");

    shared_ptr<Tracer> s1 = make_shared<Tracer>("共用");
    {
        shared_ptr<Tracer> s2 = s1;  // 兩個 shared_ptr 指向同一個物件
        cout << s1.use_count() << '\n';
    }
    cout << s1.use_count() << '\n';
    s1.reset();                      // 最後一個擁有者放手，物件銷毀
    cout << "main 結束" << '\n';
    return 0;
}
\end{example}

\iosample{}{
建立 u \\
makeTracer 結束 \\
建立 共用 \\
2 \\
1 \\
銷毀 共用 \\
main 結束 \\
銷毀 u \\
}

\begin{itemize}
\item \texttt{makeTracer} 結束時，區域變數 \texttt{p} 銷毀了，物件卻沒有
銷毀：回傳時所有權已經交給 \texttt{main} 的 \texttt{u}。 物件最後在
\texttt{main} 結束、\texttt{u} 銷毀時才跟著銷毀。
\item \texttt{s2 = s1} 之後有兩個擁有者；內層區塊結束，\texttt{s2} 銷毀，
剩一個。 \texttt{s1.reset()} 讓 \texttt{s1} 放手，擁有者變成 $0$，物件
立刻銷毀，所以 \texttt{銷毀 共用} 出現在 \texttt{main 結束} 之前。
\end{itemize}

怎麼選：先用一般變數、\texttt{string}、\texttt{vector}，大部分情況這樣
就夠了；需要 heap 上的單一物件，用 \texttt{unique\_ptr}；好幾個地方共同擁
有一個物件、不確定誰最後用完時，才用 \texttt{shared\_ptr}，它要多維護一
個計數，成本比 \texttt{unique\_ptr} 高。 照這個順序選，一般的程式碼裡就
不需要寫 \texttt{new} 和 \texttt{delete}。

\subsection{常見錯誤}

\begin{itemize}
\item \textbf{複製 \texttt{unique\_ptr}}：\texttt{unique\_ptr<int> q =
p;} 編譯失敗，訊息的開頭是 \texttt{use of deleted function
'std::unique\_ptr<\_Tp, \_Dp>::unique\_ptr(const
std::unique\_ptr<\_Tp, \_Dp>\&)'}，意思是 \texttt{unique\_ptr} 不能複
製。 要把所有權交給
\texttt{q}，寫 \texttt{unique\_ptr<int> q = std::move(p);}，下一節說明。
\item \textbf{自己 \texttt{delete} 智慧指標管的物件}：\texttt{p.get()} 會
回傳裡面的原始指標，給只收一般指標的函式用。 對它 \texttt{delete}，之後
\texttt{p} 銷毀時又 \texttt{delete} 一次，就是雙重釋放。
\end{itemize}

\section{複製與移動}
\label{sec:copy-move}

\texttt{vector<int> b = a;} 是\textbf{複製}（copy）：\texttt{b} 配置一塊
自己的記憶體，把 \texttt{a} 的元素逐一複製過去。 之後兩者各自獨立，改一
個不影響另一個。 代價跟元素個數成正比，\texttt{a} 有一百萬個元素，就要
配置約 $4$ MB 的記憶體、複製一百萬次。 這是第~2、3~章一直建議用
\texttt{const} 參考傳 \texttt{vector} 的原因。

如果 \texttt{a} 之後不會再用，可以不複製，直接把 \texttt{a} 的那塊記憶
體交給 \texttt{b}。 這叫\textbf{移動}（move），寫成
\texttt{std::move(a)}（在 \texttt{<utility>} 裡，習慣上連 \texttt{std::} 一起寫）：

\begin{example}[label=ex:copy-move]{cpp}{(copy\_move.cpp) 複製與移動 vector}
#include <iostream>
#include <utility>
#include <vector>
using namespace std;

int main() {
    vector<int> a(1000000, 7);
    vector<int> b = a;               // 複製：b 有自己的一百萬個 int
    b[0] = 1;
    cout << a[0] << ' ' << b[0] << '\n';

    const int *before = a.data();    // a 的元素放在哪裡
    vector<int> c = std::move(a);    // 移動：c 接手 a 的記憶體
    cout << c.size() << ' ' << (c.data() == before) << '\n';
    return 0;
}
\end{example}

\iosample{}{
7 1 \\
1000000 1 \\
}

\begin{itemize}
\item \texttt{v.data()} 回傳元素所在記憶體的位址。 移動之後
\texttt{c.data()} 跟原本的 \texttt{a.data()} 相同，表示 \texttt{c} 拿走
的就是 \texttt{a} 那塊記憶體，一個元素都沒有複製。 不管有幾個元素，移動
一個 \texttt{vector} 只要搬幾個指標。
\item \texttt{std::move} 本身不搬任何東西，只是把 \texttt{a} 標成「可以
被搬走」；接手記憶體的是建構 \texttt{c} 的那一步。 \texttt{string}、
\texttt{vector}、\texttt{unique\_ptr} 都會看這個標記。 第~5~章會看到，
成員是 \texttt{vector} 或 \texttt{unique\_ptr} 的類別，不用另外寫什麼就能
移動。
\item 被移走的物件（moved-from）處於\textbf{有效但未指定}（valid but
unspecified）的狀態：它還是一個合法的物件，可以重新指定、可以正常銷毀，
但內容沒有保證，除非該型別另外說明（\texttt{unique\_ptr} 就保證被移走後
是 \texttt{nullptr}）。 所以範例移動之後就不再讀 \texttt{a}；要再用，先重新指定，例如
\texttt{a = \{1, 2, 3\};}。 \texttt{string} 也一樣：
\texttt{names.push\_back(std::move(s));} 把 \texttt{s} 的字元搬進
\texttt{vector}，之後不要再讀 \texttt{s}。
\end{itemize}

很多時候移動是自動發生的，不必寫 \texttt{std::move}。 第~2~章說過，
\texttt{a + b} 這種算完就丟的值是 rvalue；暫時物件反正馬上要銷毀，編譯器
會直接移動它。 \texttt{names.push\_back(first + last)} 就會移動那個相加
出來的暫時字串。 函式回傳區域變數時也一樣，編譯器會直接把它交給呼叫者，
這就是第~3~章說回傳 \texttt{vector} 不必擔心複製的原因。

\subsection{常見錯誤}

\begin{itemize}
\item \textbf{移動之後繼續使用}：上面的 \texttt{push\_back} 之後印
\texttt{s}，編譯器不會阻止，但印出什麼沒有保證。 移走之後，除非重
新指定，否則不要再讀它。
\item \textbf{對 \texttt{const} 物件 \texttt{std::move}}：例如
\texttt{const string s = "...";} 之後寫 \texttt{string t =
std::move(s);}，編得過，也沒有
警告，但 \texttt{const} 物件不能被改，所以實際上做的是複製。 想移動的
物件不要宣告成 \texttt{const}。
\item \textbf{回傳時多寫 \texttt{std::move}}：\texttt{return
std::move(result);} 反而讓編譯器沒辦法直接把 \texttt{result} 交給呼叫者，
\texttt{g++ -Wall} 會警告 \texttt{moving a local object in a return
statement prevents copy elision}。 回傳區域變數直接寫
\texttt{return result;}。
\end{itemize}

\section{本章練習}
\label{sec:ch4-exercises}

\begin{drills}[label=drills:ch4]
\item 預測下面程式的輸出。 \texttt{T} 跟本章的 \texttt{Tracer} 一樣，只是
建立時印 \texttt{+} 和名字，銷毀時印 \texttt{-} 和名字，後面各接一個空白。
\begin{code}{cpp}{練習：建立與銷毀的順序}
void f(int k) {
    T x("x" + to_string(k));
    if (k > 1) {
        T y("y");
        return;
    }
    T z("z");
}

int main() {
    T a("a");
    f(1);
    cout << '\n';
    f(2);
    cout << '\n';
    T("tmp");
    cout << '\n';
    return 0;
}
\end{code}
\item 下面的程式用 \texttt{g++ -Wall -Wextra} 編譯沒有任何警告，但它有錯。
錯在哪裡？ 用 ASan 執行會看到什麼？ 怎麼改？
\begin{code}{cpp}{練習：找錯}
#include <cstdio>
#include <string>
using namespace std;

const char *greeting(const string &name) {
    string s = "Hello, " + name + "!";
    return s.c_str();
}

int main() {
    printf("%s\n", greeting("Barbara Liskov"));
    return 0;
}
\end{code}
\item 下面的程式要輸出陣列裡的最大值。 它確實印出 \texttt{9}，但接著在本書測試的環境下
印出 \texttt{free(): double free detected in tcache 2} 並異常結束。 裡
面有三個跟 \texttt{new}／\texttt{delete} 有關的錯誤。 找出來，再用 \texttt{vector}
改寫，讓程式裡沒有任何 \texttt{new} 和 \texttt{delete}。
\begin{code}{cpp}{練習：三個錯誤}
#include <iostream>
using namespace std;

int main() {
    int *a = new int[4]{3, 9, 2, 9};
    int *best = new int(a[0]);
    for (int i = 1; i < 4; i++) {
        if (a[i] > *best) {
            best = new int(a[i]);
        }
    }
    cout << *best << '\n';
    delete a;
    delete best;
    delete best;
    return 0;
}
\end{code}
\item 預測下面程式的輸出。 如果在最後再印一次 \texttt{b.size()}，印出的值
有保證嗎？
\begin{code}{cpp}{練習：複製、移動與智慧指標}
#include <iostream>
#include <memory>
#include <string>
#include <utility>
using namespace std;

int main() {
    string a = "abc";
    string b = a;
    string c = std::move(b);
    unique_ptr<int> p = make_unique<int>(8);
    unique_ptr<int> q = std::move(p);
    cout << a.size() << ' ' << c.size() << '\n';
    cout << (p == nullptr) << ' ' << *q << '\n';
    return 0;
}
\end{code}
\end{drills}

\addcontentsline{toc}{section}{練習解答}
\section*{練習解答}

\begin{enumerate}
\item 輸出四行：\texttt{+a +x1 +z -z -x1}、\texttt{+x2 +y -y -x2}、
\texttt{+tmp -tmp}、\texttt{-a}（每個名字後面都有一個空白；\texttt{-a} 是
\texttt{main} 結束後才印的，後面沒有換行）。 \texttt{f(1)} 依序建立 \texttt{x1}、\texttt{z}，結束時反過來銷毀。
\texttt{f(2)} 從 \texttt{if} 裡提早 \texttt{return}：先離開 \texttt{if}
的區塊，銷毀 \texttt{y}，再離開函式，銷毀 \texttt{x2}；\texttt{z} 根本沒有
建立。 \texttt{tmp} 在同一個敘述裡建立又銷毀；\texttt{a} 在最後的換行之後才銷毀。
\item \texttt{s} 是 \texttt{greeting} 的區域變數，函式結束時銷毀，
\texttt{c\_str()} 回傳的指標指向它已經釋放的記憶體，\texttt{printf} 讀的
是懸置指標。 用 ASan 執行會看到 \texttt{heap-use-after-free}。 改法是讓
\texttt{greeting} 回傳 \texttt{string}，呼叫端寫
\texttt{printf("\%s\textbackslash n", greeting("Barbara
Liskov").c\_str());}，輸出 \texttt{Hello, Barbara Liskov!}。
暫時字串活到分號，\texttt{printf} 用完時它還在。
\item 三個錯誤：(1) 迴圈裡每找到更大的數就 \texttt{new} 一個新的
\texttt{int}，舊的那個沒有 \texttt{delete}，洩漏；(2) \texttt{a} 是
\texttt{new[]} 配置的，要用 \texttt{delete[]}，寫成 \texttt{delete} 是
未定義行為；(3) \texttt{best} 被 \texttt{delete} 兩次。 改寫時陣列換成
\texttt{vector<int> a = \{3, 9, 2, 9\};}，最大值用一個 \texttt{int best}
記就夠了，迴圈裡寫 \texttt{if (a[i] > best) best = a[i];}，整支程式不需要
任何 \texttt{new} 或 \texttt{delete}。
\item 輸出兩行：\texttt{3 3} 和 \texttt{1 8}。 \texttt{a} 沒有被動過，還是
$3$ 個字元；\texttt{c} 拿到 \texttt{b} 從 \texttt{a} 複製來的 \texttt{"abc"}；
\texttt{p} 被移走之後保證是 \texttt{nullptr}，印 \texttt{1}；\texttt{*q}
是 $8$。 \texttt{b} 是被移走的 \texttt{string}，處於有效但未指定的狀態，
所以 \texttt{b.size()} 沒有保證；要用 \texttt{b}，先重新指定或呼叫
\texttt{b.clear()}。
\end{enumerate}
